\begin{papercard}{tangari2021}
\cf{Stages} \stagetag{A} \stagetag{D} \stagetag{R}. A census-scale cross-sectional study of 20{,}991 Android mHealth apps (Australian Google Play). It combines static code analysis, automated traffic capture, machine-learning policy analysis and review mining.

\cf{Aim} Investigate whether and what user data mHealth apps collect, characterise the privacy conduct of all available Google Play mHealth apps, and gauge the risks \pc{tangari2021}{1}{To investigate whether and what user data are}. Outcomes: data collection in code, transmissions in traffic, recipients, adverts/trackers, policy audit with compliance, and complaints in negative reviews \pc{tangari2021}{1}{Primary outcomes were characterisation of the data}.

\cf{Sampling frame (distinctive)} Google Play does not list all mHealth apps, and its search does not show all available apps. A crawler therefore started from the top 100 Medical and Health \& Fitness apps and \emph{snowballed through ``similar apps''}, crawling more than 1.7 million apps between 1 October and 15 November 2019 \pc{tangari2021}{2}{Google Play does not provide a complete list of}. Apps in the two categories were kept: 20{,}991 in total (75.7\,\% free, 15.4\,\% paid, 8.9\,\% geoblocked). A random set of 8{,}468 popular non-health apps served as a baseline comparator \pc{tangari2021}{2}{we used the crawler to sample a random set}. Most apps were not Australia-specific, so the frame is argued to proxy the worldwide store \pc{tangari2021}{2}{From this information we}.

\cf{Policy capture} For all 15{,}893 free apps the public policy was extracted, typically via the store link. A broken link or an empty page counted as \emph{no privacy policy} \pc{tangari2021}{2}{If the link was broken or directed users to a page}. This is an explicit operational rule for absence.

\cf{Policy analysis} Manual annotation was infeasible at this scale. A supervised classifier predicted disclosure of personal-data types, trained on APP-350 (350 mobile-app policies annotated by legal experts). Reported performance: accuracy $>$97\,\%, mean precision 87\,\%, mean recall 77\,\% \pc{tangari2021}{4}{To overcome this, we analysed the automatic}.

\cf{Behavioural analyses} Static analysis: apktool decoding, dictionary search over 338 third-party libraries, and Android API calls that access personal data, kept only if the needed permission was declared (a superset of actual collection) \pc{tangari2021}{2}{we processed the downloaded app packages using}. Traffic: a testbed with a phone, a Wi-Fi access point and interception, running about 35 automated UI actions in a 180-second session per app \pc{tangari2021}{4}{we performed 35 different activities}. Adverts and trackers were identified with EasyList/EasyPrivacy. PII in traffic was found with a classifier (97\,\% accuracy). The authors call this a lower bound, because automated testing may not trigger all views \pc{tangari2021}{4}{The result only includes data collection}. Reviews: 2.1 million reviews of 6{,}938 apps, split by stars, with 41 keywords mapped to six complaint categories \pc{tangari2021}{6}{Supplementary appendix F lists 41}.

\cf{Policy-vs-behaviour classification} Each observed transmission was classed as \emph{complying} (disclosed in the policy), \emph{violating} (the app has a policy but the practice is not disclosed) or \emph{no privacy policy} \pc{tangari2021}{6}{each data transmission was classed as complying}.

\cf{Key findings} 88\,\% of apps had code able to collect personal data, and 3.9\,\% transmitted it in testing. 87.5\,\% of collection operations stemmed from third-party libraries, and 50 third parties accounted for about 68\,\% \pc{tangari2021}{5}{The results show a predominant role of}. 28.1\,\% had no valid policy text. Policy presence rose with installs, to 94.4\,\% for apps with $\geq$1 million installs \pc{tangari2021}{6}{A positive correlation was also found between an}. 47\,\% of transmissions complied with the policy, 24\,\% violated it and 29\,\% came from apps without a policy (Table 5). 23\,\% of transmissions were unencrypted. Privacy complaints made up 1.3\,\% of reviews.

\cf{Limitations stated} ML false positives, randomised rather than scripted interaction that misses parts of apps, and free apps only \pc{tangari2021}{9}{To scale up the study and cope with a large number}.

\cf{Recommendations} Clinicians should check permissions and policies. Stores should vet policy links, since links were often broken or led to empty pages \pc{tangari2021}{10}{we observed that the links to privacy}.

\cf{Use for this thesis} Mainly a \emph{contrast} case: it shows what a communication-only study cannot claim, namely observed transmissions. Directly reusable: the broken/empty-link rule for policy absence, the comply/violate/no-policy trichotomy as an analogy for coding \emph{disclosure gaps} without claiming behaviour, the evidence that popularity correlates with policy presence (a sampling-bias argument), and the ``similar apps'' crawl as a frame-building idea. Out of scope: traffic, static analysis and ML policy classification.
\end{papercard}
